\TieudeT{PHONG TỎA VẬT LÍ 12 - VẬT LÍ NHIỆT}
\subsubsection*{PHẦN I. Thí sinh trả lời câu hỏi từ câu 1 đến câu 18. Mỗi câu hỏi thí sinh chỉ chọn một phương án}
\setcounter{ex}{0}
\begin{ex}%chuan
	Trường hợp nào dưới đây làm biến đổi nội năng \textbf{không} do thực hiện công?
	\choice
	{Mài dao}
	{Đóng đinh}
	{Khuấy nước}
	{\True Nung sắt trong lò}
	\loigiai{}
\end{ex}


\begin{ex}%chuan
	Quá trình nào sau đây là quá trình tỏa nhiệt?
	\choice
	{Nóng chảy - thăng hoa}
	{Bay hơi - thăng hoa}
	{\True Ngưng kết - ngưng tụ}
	{Nóng chảy - bay hơi}
	\loigiai{}
\end{ex}

\begin{ex}%chuan
	Phát biểu nào sau đây là đúng?
	\choice
	{Nội năng của một vật phụ thuộc vào nhiệt độ của vật}
	{Nội năng của một vật phụ thuộc thể tích của vật}
	{\True Nội năng của một vật phụ thuộc thể tích và nhiệt độ của vật}
	{Nội năng của một vật không phụ thuộc thể tích và nhiệt độ của vật}
	\loigiai{}
\end{ex}

\begin{ex}%chuan
	Phát biểu nào sau đây là \textbf{sai} khi nói về cấu tạo chất?
	\choice
	{Các phân tử luôn luôn chuyển động không ngừng}
	{Các phân tử chuyển động càng nhanh thì nhiệt độ của vật càng cao và ngược lại}
	{\True Các phân tử luôn luôn đứng yên và chỉ chuyển động khi nhiệt độ của vật càng cao}
	{Các chất được cấu tạo từ các hạt riêng biệt là phân tử}
	\loigiai{}
\end{ex}

\begin{ex}%chuan
	Trong sự dẫn nhiệt, nhiệt chỉ có thể tự truyền
	\choice
	{từ vật có khối lượng lớn hơn sang vật có khối lượng nhỏ hơn}
	{từ vật có thể tích lớn hơn sang vật có thể tích nhỏ hơn}
	{\True từ vật có nhiệt độ cao hơn sang vật có nhiệt độ thấp hơn}
	{từ vật có nhiệt năng lớn hơn sang vật có nhiệt năng nhỏ hơn}
	\loigiai{}
\end{ex}
\begin{ex}%chuan
	Đặc điểm nào sau đây là của thể lỏng?
	\choice
	{Khoảng cách giữa các phân tử rất lớn so với kích thước của chúng}
	{\True Lực tương tác phân tử yếu hơn lực tương tác phân tử ở thể rắn}
	{Không có thể tích và hình dạng riêng xác định}
	{Các phân tử dao động xung quanh vị trí cân bằng xác định}
	\loigiai{}	
\end{ex}

\begin{ex}%chuan
	Không thể dùng nhiệt kế rượu để đo nhiệt độ của hơi nước đang sôi vì
	\choice
	{rượu sôi ở nhiệt độ cao hơn $100^{\circ}C$}
	{\True rượu sôi ở nhiệt độ thấp hơn $100^{\circ}C$}
	{rượu đông đặc ở nhiệt độ cao hơn $100^{\circ}C$}
	{rượu đông đặc ở nhiệt độ thấp hơn $0^{\circ}C$}
	\loigiai{}
\end{ex}

\begin{ex}%chuan \nguon{4}
	Trong chuyển động nhiệt, các phân tử chất lỏng
	\choice
	{chuyển động hỗn loạn quanh vị trí cân bằng}
	{chuyển động hỗn loạn quanh vị trí cân bằng xác định}
	{chuyển động hỗn loạn}
	{\True dao động quanh vị trí cân bằng nhưng những vị trí này không cố định mà di chuyển}
	\loigiai{}
\end{ex}

\begin{ex}
	Cho hiện tượng sau: Tơ nhện được hình thành từ một loại protein dạng lỏng trong cơ thể nhện. Khi làm tơ, nhện nhả protein đó ra khỏi cơ thể, protein đó sẽ chuyển thành tơ nhện. Hiện tượng trên thể hiện sự chuyển thể của protein là
	\choice
	{sự bay hơi}
	{sự nóng chảy}
	{\True sự đông đặc}
	{sự ngưng tụ}
	\loigiai{}
\end{ex}

\begin{ex}
	Người ta thực hiện công $40\ J$ để nén khí trong một xilanh. Khí truyền ra môi trường xung quanh một nhiệt lượng $10\ J$. Độ biến thiên nội năng của khí là
	\choice
	{\True $30\ J$}
	{$50\ J$}
	{$20\ J$}
	{$60\ J$}
	\loigiai{}
\end{ex}

\begin{ex}
	Biết nhiệt nóng chảy riêng của bạc là $0,88.10^5\ J/kg$. Nhiệt lượng cần cung cấp cho $525\ g$ bạc nóng chảy ở nhiệt độ nóng chảy là
	\choice
	{$462.10^5\ J$}
	{$0,462\ J$}
	{\True $4, 62.10^4\ J$}
	{$1, 68.10^5\ J$}	
	\loigiai{}
\end{ex}

\begin{ex}
	Biết nhiệt nóng chảy riêng của nhôm là $3,9.10^5\ J/kg$. Để hóa lỏng hoàn toàn miếng nhôm đó ở nhiệt độ $659^{\circ}C$ cần nhiệt lượng là $20.10^5\ J$. . Khối lượng của miếng nhôm là
	\choice
	{$0,195\ g$}
	{$5,13\ g$}
	{$0,195\ kg$}
	{\True $5,13\ kg$}
	\loigiai{}
\end{ex}

\begin{ex}
	Biết nhiệt dung riêng của gỗ là $c = 1240\  J.kg^{-1}.K^{-1}$. Khi $100\ g$ gỗ giảm nhiệt độ đi $1\ K$ thì nó
	\choice
	{cần nhận nhiệt lượng $124\ J$ từ môi trường bên ngoài}
	{\True giải phóng một năng lượng bằng $124\ J$ ra môi trường bên ngoài}
	{giải phóng một năng lượng bằng $12,4\ J$ ra môi trường bên ngoài}
	{cần nhận nhiệt lượng $1240\ J$ từ môi trường bên ngoài}
	\loigiai{}
\end{ex}

\begin{ex}
	 Cho biết nước đá có nhiệt nóng chảy riêng là $3,4.10^5\ J/kg$ và nhiệt dung riêng là $2,09.10^3\ J/kg.K$, nước đá nóng chảy ở $0^{\circ}C$. Nhiệt lượng $Q$ cần cung cấp để làm nóng chảy cục nước đá khối lượng $100\ g$ đang có nhiệt độ $-25^{\circ}C$ là
	\choice
	{$34000\ J$}
	{$5225\ J$}
	{\True $39225\ J$}
	{$28775\ J$}	
	\loigiai{}
\end{ex}

\begin{ex}
	Một máy khoan công suất $1\ kW$ dùng để khoan một cái lỗ trên tấm nhôm khối lượng $0,8\ kg$. Biết rằng $50\%$ điện năng được sử dụng để làm nóng máy và tỏa nhiệt ra môi trường xung quanh. Biết nhiệt dung riêng của nhôm là $880\ J/kg.^\circ C$. Độ tăng nhiệt độ của khối nhôm sau $2,5$ phút là
	\choice
	{$66,5\ ^\circ C$}
	{$106,5\ ^\circ C$}
	{$100,5\ ^\circ C$}
	{$90,5\ ^\circ C$}
	\loigiai{
		\begin{itemize}
			\item  Ta có: $Q = m.c.\Delta t \Rightarrow \Delta t = \dfrac{Q}{mc} = \dfrac{0,5.\mathcal{P}.\tau}{mc} = \dfrac{0,5.10^3.150}{0,8.880} \approx 106,5^\circ C$.
		\end{itemize}
	}
\end{ex}

\begin{ex}
	Hai miếng nhôm và chì có cùng khối lượng rơi từ cùng một độ cao xuống sàn nhà. Nhiệt dung riêng của nhôm là $880\ J/kg.K$, của chì là $130\ J/kg.K$. Tỉ số độ tăng nhiệt độ của hai miếng kim loại trên khi chúng va chạm với sàn nhà nếu coi toàn bộ cơ năng của vật khi rơi đều dùng để làm nóng vật là
	\choice
	{\True $6,77$ lần}
	{$6$ lần}
	{$5,77$ lần}
	{$5$ lần}
	\loigiai{}
\end{ex}

\begin{ex}
	Súng là một dạng của động cơ nhiệt và hoạt động như một động cơ đốt trong mà không có pit-tông hoạt động theo chu trình, thay vào đó, nó tham gia vào một quá trình giãn nở nhiệt đơn lẻ trong đó đầu đạn được phóng đi. Với khối lượng thân súng làm bằng sắt là $1,8\ kg$, súng này phóng một viên đạn có khối lượng $2,4\ g$ với vận tốc $320\ m/s$ và có hiệu suất $1,1\%$. Giả sử rằng thân súng hấp thụ toàn bộ nhiệt năng thừa. Bỏ qua sự mất mát nhiệt vào môi trường xung quanh. Biết rằng nhiệt dung riêng của sắt là $448\ J/kg.K$. Nhiệt độ tăng lên của súng sau khi bắn một phát đạn là
	\choice
	{\True $13,7\ ^\circ C$}
	{$11,7\ ^\circ C$}
	{$12,0\ ^\circ C$}
	{$10,0\ ^\circ C$}
	\loigiai{
	\begin{itemize}
		\item Độ biến thiên động năng: $W_{\text{đ}} = \dfrac{1}{2}mv^2 = \dfrac{1}{2}.2,4.10^{-3}.320^2 = 122,88\ J$.
		\item Ta có: $H = \dfrac{W_{\text{đ}}}{W_{\text{đ}}+Q} \Rightarrow Q =\approx 11048\ J$.
		\item $Q = m_s.c.\Delta t \Rightarrow \Delta t = \dfrac{Q}{m_s.c} = \dfrac{11048}{1,8.448} \approx 13,7^\circ C$.
	\end{itemize}
	}
\end{ex}

\begin{ex}
	Một dòng điện cường độ $I = 10\ A$ chạy qua dây dẫn bằng đồng có điện trở $R = 2\ \Omega$ và có khối lượng $200\ g$. Biết rằng nhiệt dung riêng của đồng là $c = 880\ J/kg.K$. Cho rằng toàn bộ nhiệt lượng tỏa ra trên dây dẫn được dây dẫn hấp thụ và tăng nhiệt độ. Độ tăng nhiệt độ của dây dẫn sau khoảng thời gian $5$ phút dòng điện chạy qua nó là
	\choice
	{$111,0\ ^\circ C$}
	{$141,5\ ^\circ C$}
	{\True $340,9\ ^\circ C$}
	{$242,6\ ^\circ C$}
	\loigiai{
	\begin{itemize}
		\item $Q = m.c.\Delta t \Rightarrow I^2.R.\tau = m.c.\Delta t \Rightarrow \Delta t = \dfrac{I^2.R.\tau}{m.c} = \dfrac{10^2.2.5.60}{0,2.880} \approx 340,9^\circ C$.
	\end{itemize}
	}
\end{ex}
\subsubsection*{PHẦN 2. Thí sinh trả lời từ câu 1 đến câu 4. Trong mỗi ý a), b), c), d) ở mỗi câu, thí sinh chọn đúng hoặc sai.}
\setcounter{ex}{0}
\begin{ex}
	\immini[thm]{Khi làm thí nghiệm đun nóng một chất. Kết quả thí nghiệm, một học sinh vẽ được biểu diễn sự thay đổi nhiệt độ của chất đó theo thời gian như hình bên.
		}
		{
		\begin{tikzpicture}[draw=Mapcolor, very thick, >=latex']
			\draw [->] (0,0) -- (6,0) node [above] {Thời gian};
			\draw [->] (0,0) -- (0,3) node [right] {Nhiệt độ $(^\circ C)$};
			\draw (0,0) node [left]{$O$} -- (0.25,0.5) circle (1pt) node [right] {$A$} -- (0.5,1) -- (1.5,1) circle (1pt) node [above] {$B$} -- (2.5,1) -- (3,1.75) circle (1pt) node [left] {$C$} -- (3.5,2.5) -- (4.5,2.5) circle (1pt) node [above] {$D$} -- (5.5,2.5) ;
			\draw [dashed]
			(0.25,0.5) -- (0.25,0) node [below] {$t_1$}
			(1.5,1) -- (1.5,0) node [below] {$t_2$}
			(3,1.75) -- (3,0) node [below] {$t_3$}
			(4.5,2.5) -- (4.5,0) node [below] {$t_4$}
			(1.5,1) -- (0,1) node [left] {$17$}
			(4.5,2.5) -- (0,2.5) node [left] {$115$};
		\end{tikzpicture}
		}
	\choiceTFt
	{Tại thời điểm $t_2$, chất ở thể lỏng}
	{\True Nhiệt độ nóng chảy của chất là $17\ ^{\circ}C$}
	{Tại thời điểm $t_3$, chất ở thể rắn và thể lỏng}
	{\True Nhiệt độ sôi của chất này là $115\ ^{\circ}C$}
	\loigiai{
		\begin{itemchoice}
			\itemch 
			\itemch
			\itemch 
			\itemch 
		\end{itemchoice}
	}
\end{ex}

\begin{ex}
	Để xác định nhiệt nóng chảy riêng của nước đá, có thể tiến hành thí nghiệm theo sơ đồ nguyên lí như hình dưới đây. Dòng điện làm nóng dây điện trở trong một nhiệt lượng kế và làm nước đá nóng chảy. Lượng nuớc thu được sau khi toàn bộ nước đá nóng chảy được đem đi cân thì thấy nó có khối lượng $15\ g$. Công suất điện tiêu thụ được xác định bằng oát kế là $24\ W$. Thời gian đun được xác định bằng đồng hồ là $180\ s$. Bỏ qua sự trao đổi nhiệt với môi trường xung quanh.
	\begin{center}
		\begin{tikzpicture}[draw=Mapcolor, very thick, scale = 1, line join=round, line cap=round, >=stealth,line width=1.5]
			%\draw[dashed,line width=0.5] (0,0) grid (8,6);
			\draw[rounded corners=5] (0,0) rectangle (2,1.5);
			\draw[rounded corners=2] (0.5,1) rectangle (1.5,1.2);
			\draw (0.4,0.6) circle (3pt) (1.6,0.6) circle (3pt);
			\node at (0.7,0.3) {+};
			\node at (1.4,0.3) {$\boldsymbol{-}$};
			\node at (1,2) {Biến áp nguồn};
			\draw[rounded corners=5] (2.5,0.4) rectangle (3.5,1.2);
			\node at (3,1.7) {Oát kế};
			\foreach \x/\y in {0.1/0.1,0.6/0.3,1.5/0.9,0.4/1,0/0.5,1/0.2,1.4/0.3,-0.1/1.1,0.7,1.1}{\draw[ball color=Mapcolor] (5.8+\x,0.3+\y) circle (5pt);}
			\draw[rounded corners=5] (5,3) -- (5.5,0) -- (7.5,0) -- (8,3);
			\draw[rounded corners=2,fill=gray] (4.8,3) rectangle (8.2,3.2);
			\draw[decorate,decoration={coil,aspect=0}]  (6,1) -- (7,1);
			\draw (6,3) -- (6,1) (7,3) -- (7,1);
			\draw[fill=black] (2.7,1) circle (2pt) (2.7,0.6) circle (2pt);
			\draw[fill=black] (3.3,1) circle (2pt) (3.3,0.6) circle (2pt);
			\draw (2,0.6) -- (2.7,0.6) (2,1) -- (2.7,1);
			\draw (3.3,0.6) -- (4.1,0.6) -- (4.1,4) -- (6,4) -- (6,3.2);
			\draw (3.3,1) -- (3.9,1) -- (3.9,4.2) -- (7,4.2) -- (7,3.2);
			\draw[rounded corners=1,fill=gray] (5.9,3.2) rectangle (6.1,3.6);
			\draw[rounded corners=1,fill=gray] (6.9,3.2) rectangle (7.1,3.6);
			\draw (6.5,0.9) -- (8,0.2) node [right] {Dây điện trở};
			\draw (7.5,3.2) -- (8,4) node [right] {Nhiệt lượng kế};
		\end{tikzpicture}
	\end{center}
	\choiceTFt
	{\True Ở áp suất tiêu chuẩn, nước đá tinh khiết nóng chảy ở nhiệt độ $0{}\ ^{\circ} C$}
	{Khi nóng chảy, thể tích nước tăng}
	{\True Nhiệt lượng nước đá nhận được từ dây điện trở là 4320 J}
	{Nhiệt nóng chảy riêng $\lambda$ của nước đá đo được là $3,3 . 10^5 \ J/kg$}
	\loigiai{
		\begin{itemchoice}
			\itemch 
			\itemch 
			\itemch 
			\itemch 
		\end{itemchoice}
	}
\end{ex}

\begin{ex}
	
	\immini[thm]
	{Để tạo ra các loại rượu truyền thống đặc trưng của Việt Nam, các cơ sở sản suất rượu đã thực hiện nhiều công đoạn. Trong các công đoạn đó có một công đoạn rất quan trọng là quá trình chưng cất rượu. Quá trình chưng cất được thể hiện đơn giản bằng sơ đồ như hình bên.}
	{\begin{tikzpicture}[draw=Mapcolor, very thick, >=latex']
			%\draw [dotted] (-3,0) grid (4,6);
			\fill[bottom color=cyan, top color=white] (-1.5,0) rectangle (1.5,2.5);
			\draw[line width=1.5pt] (-0.25,4) -- (-1.5,3.5) -- (-1.5,0) -- (1.5,0) -- (1.5,3.5) -- (0.25,4);
			\draw [double distance=0.5cm, line width=1.5pt, rounded corners] (0,4) -- (0,5) -- (4,5) -- (4,4) ;
			\draw (3,0) rectangle (5,4);
			\foreach \x in {3.1,3.5,3.9,4.3,4.7}{
				\draw [fill = blue!50] (\x,0) rectangle (\x+0.2,4) ;
			}
			\draw
			node[text width=2cm, align=center]at (0,1) {Hỗn hợp nhiên liệu hóa lỏng}
			node [below] at (0,0) {$A$}
			node [below] at (4,0) {$B$}
			node [above] at (2,4) {Ống dẫn hơi rượu}
			node [above,rotate = -90] at (5,2) {Nước làm mát};
			\draw [->, red] (0,2.5) -- (0,3.5);
			\draw [->, red] (-1,3) -- (-0.2,3.8);
			\draw [->, red] (1,3) -- (0.2,3.8);
	\end{tikzpicture}}
	\choiceTFt
	{\True Ở bồn A hỗn hợp nguyên liệu lỏng được cung cấp nhiệt lượng để tạo hơi rượu}
	{\True Ở bồn B diễn ra quá trình hơi rượu ngưng tụ thành rượu dưới dạng lỏng}
	{Nước làm mát ở bồn B có tác dụng cung cấp nhiệt lượng cho quá trình ngưng tụ của hơi rượu}
	{\True Thực ra trong hỗn hợp hơi vừa có hơi rượu vừa có hơi nước}
	
	\loigiai{
		\begin{itemchoice}
			\itemch 
			\itemch 
			\itemch Vì nước thu nhiệt chứ không cấp nhiệt
			\itemch Vì ở bồn B khi được cung cấp nhiệt lượng thì nước cũng hóa hơi cùng với rượu nhưng ít hơn do nhiệt độ sôi của nước cao hơn.
		\end{itemchoice}
	}
\end{ex}

\begin{ex}
	Một khối gỗ có khối lượng $M = 200\ g$ được đặt trên một mặt phẳng nằm ngang nhẵn. Một viên đạn chì khối lượng $m = 20\ g$ (nhiệt dung riêng của chì là $126\ J/kg.K$) bay theo phương ngang với tốc độ $v_0 = 500\ m/s$ được bắn vào khối gỗ. Sau khi đạn xuyên qua, khối gỗ có tốc độ $v_g = 30\ m/s$. Biết 42\% cơ năng mất đi của quá trình này chuyển hóa thành nội năng của viên đạn (làm nóng viên đạn).
	\choiceTFt
	{Tốc độ của viên đạn sau khi xuyên qua khối gỗ là $20\ m/s$}
	{\True Cơ năng của hệ (đạn + khối gỗ) ban đầu là $2500\ J$}
	{Tổng nhiệt lượng sinh ra trong quá trình này là $1560\ J$}
	{\True Nhiệt độ của viên đạn tăng lên là $\Delta t = 335^\circ C$}
	\loigiai{
		\begin{itemchoice}
			\itemch Bảo toàn động lượng: $mv_0 = MV + mv \Rightarrow 0,02.500 = 0,2.30+0,02.v \Rightarrow v = 200 m/s$.
			\itemch $W = \dfrac{1}{2}mv_0^2 = 0,5.0,02.500^2 = 2500\ J$.
			\itemch $\Delta W = W - W_0 = \dfrac{1}{2}mv_0^2 - \left( \dfrac{1}{2}MV^2 + \dfrac{1}{2}mv_s^2\right)  = 2500 - \left(0,5.0,2.30^2+0,5.0,02.200^2\right) =2010\ J$.
			\itemch $Q = 0,42\Delta W \Rightarrow m.c.\Delta t = 0,42.\Delta W \Rightarrow \Delta t = \dfrac{0,42.\Delta W}{mc} = \dfrac{0,42.2010}{0,02.126} = 335^\circ C$.
		\end{itemchoice}
	}
\end{ex}
\subsubsection*{PHẦN 3. Thí sinh trả lời từ câu 1 đến câu 6.}
\setcounter{ex}{0}
\begin{ex}% Câu 1
	Người ta cọ xát hai vật với nhau, nhiệt dung của hai vật bằng nhau và bằng $800\ J/K$. Sau $1$ phút người ta thấy nhiệt độ của mỗi vật tăng thêm $30\ K$. Công suất trung bình của việc cọ xát bằng bao nhiêu $W$?
	\SA[oly]{$800$}
	\loigiai{Toàn bộ công cọ xát chuyển thành nhiệt của hai vật: $$A = Q \Rightarrow \mathcal{P}\tau = 2mc\Delta t \Rightarrow \mathcal{P}.60 =2.800.30 \Rightarrow \mathcal{P} = 800\ W $$.}
\end{ex}

\begin{ex}% Câu 2
	Có ba bình nước giống nhau, mỗi bình chứa $20\ g$ nước ở cùng nhiệt độ. Người ta thả vào mỗi bình một cục nước đá có khối lượng khác nhau nhưng có cùng nhiệt độ. Bình 1 được thả cục nước đá có khối lượng $10\ g$, khi cân bằng nhiệt, khối lượng nước đá trong bình là $9\ g$. Bình 2 được thả cục nước đá có khối lượng $20\ g$, khi cân bằng nhiệt, khối lượng nước đá trong bình không đổi. Bình 3 được thả cục nước đá có khối lượng $40\ g$, khi có cân bằng nhiệt, khối lượng nước đá trong bình là bao nhiêu gam?
	\SA[oly]{$42$}
	\loigiai{
	\begin{itemize}
		\item Bình 1 có $\Delta m_1 = 10-9 = 1\ g$ tan thành nước nên nhiệt độ cân bằng là $0^\circ C$:
		$$ m_nc_nt_n = m_{\text{đ1}}c_{\text{đ}}\left(0-t_{\text{đ}} \right) + \Delta m_1.\lambda \Rightarrow 20c_nt_n = 10c_{\text{đ}}\left(0-t_{\text{đ}} \right) +\lambda (1)$$ 
		\item Bình 2 có khối lượng nước đá không đổi nên nhiệt độ cân bằng là $0^\circ C$:
		$$ m_nc_nt_n = m_{\text{đ2}}c_{\text{đ}}\left(0-t_{\text{đ}} \right) \Rightarrow 20c_nt_n = 20c_{\text{đ}}\left(0-t_{\text{đ}} \right) (2)$$ 
		\item Từ (1) và (2) $\Rightarrow \lambda = 10c_{\text{đ}}\left(0-t_{\text{đ}} \right)$.
		\item Bình 3 thả cục nước đá $40\ g$ nên một phần nước sẽ đông đặc:
		$$ m_nc_nt_n + m_{\text{đđ}}\lambda = m_{\text{đ3}}c_{\text{đ}}\left( 0-t_{\text{đ}}\right) \Rightarrow 20c_nt_n +  m_{\text{đđ}}\lambda =40c_{\text{đ}}\left( 0-t_{\text{đ}}\right)  \Rightarrow m_{\text{đđ}} = 2\ g$$
	\end{itemize}
	}
\end{ex}

\begin{ex}% Câu 3
	Một nhà máy nhiệt điện tiêu thụ $0,35\ kg$ nhiên liệu cho mỗi kWh điện. Cho biết năng suất tỏa nhiệt của nhiên liệu là $42.10^6\ J/kg$. Hiệu suất thực của động cơ nhiệt dùng trong nhà máy điện bằng bao nhiêu phần trăm (làm tròn kết quả đến chữ số hàng đơn vị)?	
	\shortans[oly]{$24$}
	\loigiai{
		\begin{itemize}
			\item $H = \dfrac{A_{ci}}{Q}.100\% = \dfrac{1000.3600}{0,35.42.10^6}.100\% \approx 24,48\% \approx 24\%$.
		\end{itemize}
	}
\end{ex}

\begin{ex}% Câu 4
	Vận động viên điền kinh bị mất rất nhiều nước trong khi thi đấu. Các vận động viên thường chỉ có thể chuyển hoá khoảng $20\%$ năng lượng dự trữ trong cơ thể thành năng lượng dùng cho các hoạt động của cơ thể. Phần năng lượng còn lại chuyển thành nhiệt thải ra ngoài nhờ sự bay hơi của nước qua hô hấp và da để giữ cho nhiệt độ cơ thể không đổi. Coi nhiệt độ cơ thể của vận động viên hoàn toàn không đổi và nhiệt hoá hơi riêng của nước ở nhiệt độ của vận động viên là $2,4.10^{6}\ J/kg$. Biết khối lượng riêng của nước là $10^{3}\ kg/m^{3}$. Nếu vận động viên dùng hết $10800\ kJ$ trong cuộc thi thì có bao nhiêu lít nước đã thoát ra ngoài cơ thể?
	\shortans[oly]{$3,6$}
	\loigiai{
		\begin{itemize}
			\item $V = \dfrac{m}{D} = \dfrac{Q}{L.D} = \dfrac{(1-0,2).10800.10^3}{2,4.10^6.10^3} = 3,6.10^{-3}\ m^3 = 3,6$ lít.
		\end{itemize}
	}
\end{ex}

\begin{ex}% Câu 5
	Có hai quả cầu kim loại cùng kích thước, có động năng lần lượt là $100\ J$ và $128\ J$, chuyển động cùng phương ngang ngược chiều đến va chạm mềm trực diện với nhau. Quả cầu có động năng bé hơn có khối lượng gấp đôi quả cầu còn lại. Cho rằng, toàn bộ phần cơ năng bị giảm chuyển thành nội năng của hai quả cầu. Độ tăng nội năng của hệ hai quả cầu là bao nhiêu $J$ (làm tròn kết quả đến hàng đơn vị)?
	\shortans[oly]{$225$}
	\loigiai{
		\begin{itemize}
			\item Hai vật va chạm mềm:
			$\begin{cases}
				p&=mv\\
				W_{\text{đ}} &= \dfrac{1}{2}mv^2
			\end{cases}
			\Rightarrow p^2 = 2mW_{\text{đ}}$
		
		\item Ta có: $\overrightarrow{p_1} + \overrightarrow{p_2} = \overrightarrow{p} \Rightarrow \left| p_1 - p_2\right|  = p  \Rightarrow \left| \sqrt{2m_1W_{\text{đ1}}} - \sqrt{2m_2W_{\text{đ2}}}\right| = \sqrt{2(m_1+m_2)W_{\text{đ}}} \Rightarrow \left| \sqrt{2.100} -\sqrt{1.128}\right|  = \sqrt{(2+1)W_{\text{đ}}} \Rightarrow W_{\text{đ}} = \dfrac{8}{3}\ J$.
		\item $\Delta U = W_{\text{đ1}} + W_{\text{đ2}} - W_{\text{đ}} = 100 + 128 - \dfrac{8}{3} \approx 225,3 \approx 225\ J$.
	\end{itemize}
	}
\end{ex}

\begin{ex}% Câu 6
	Khi trời lạnh, cơ thể chúng ta dễ bị mất nhiệt lượng vào môi trường. Sự mất nhiệt lượng này sẽ dẫn đến các hệ quả như cơ thể bị run rẩy, kiệt sức. Trong trận bóng đá ngoài trời vào ngày lạnh, cầu thủ sẽ bắt đầu cảm thấy kiệt sức sau khi tiêu hao khoảng $8,1.10^5\ J$ nội năng. Biết nhiệt lượng do cơ thể một cầu thủ truyền cho môi trường là $6,5.10^5\ J$. Công mà cầu thủ này đã thực hiện có độ lớn bằng bao nhiêu $MJ$?	
	\shortans[oly]{$0,15$}
	\loigiai{
		\begin{itemize}
			\item $\Delta U = A + Q \Rightarrow -8.10^5 = A - 6,5.10^5 \Rightarrow A = -0,15.10^6 = -0,15\ MJ$.
		\end{itemize}
	}
\end{ex}
